\chapter{Całki Feynmana, obrót Wicka i mechanika kwantowa}
\label{ch:feynman-path}

\section{Dwa ewolucyjne równania i dwie rodziny operatorów}

W mechanice kwantowej funkcja falowa $\psi(t,\cdot)$ zmienia się
zgodnie z równaniem Schrödingera. Dla cząstki o masie $m>0$
i stałej Plancka $\hbar>0$ rozważmy w $L^2(\R^d)$
operator
\[
 H_0=-\frac{\hbar^2}{2m}\Delta,
 \qquad H=H_0+V,
\tag{16.1}\label{eq:QM-H}
\]
gdzie $V:\R^d\to\R$ jest rzeczywistą, ograniczoną funkcją
borelowską i $V$ oznacza operator mnożenia przez tę funkcję.
Dziedziną $H_0$ jest $H^2(\R^d)$, a $H$ jest
samiosprzężony na tej samej dziedzinie: dodanie ograniczonego
operatora samosprzężonego zachowuje samosprzężoność.
Wtedy $H\ge-\|V\|_\infty I$. Z twierdzenia spektralnego
otrzymujemy dwie rodziny operatorów:
\[
 U(t)=e^{-itH/\hbar},\quad t\in\R,\qquad
 S(s)=e^{-sH/\hbar},\quad s\ge0.
\tag{16.2}\label{eq:unitary-euclidean}
\]
Pierwsza jest \emph{grupą unitarną}, druga --- \emph{półgrupą
cieplną} (niekoniecznie kontrakcją, jeśli $V$ bywa ujemny).

\begin{stwierdzenie}[W jakim sensie operatory rozwiązują równania?]
\label{prop:unitary-generator}
Niech $f\in H^2(\R^d)$. Wówczas $u(t)=U(t)f$ jest
różniczkowalne jako funkcja o wartościach w $L^2(\R^d)$
i spełnia $i\hbar u'(t)=Hu(t)$, $u(0)=f$.
Podobnie $v(s)=S(s)f$ spełnia $\hbar v'(s)=-Hv(s)$.
Dla dowolnego $f\in L^2$ operatorowe wzory
\eqref{eq:unitary-euclidean} nadal mają sens, ale
pochodne w $L^2$ nie muszą istnieć.
\end{stwierdzenie}
\begin{proof}
Dla miary spektralnej $E_H$ mamy
$U(t)f=\int_{\R}e^{-it\lambda/\hbar}\,\dd E_H(\lambda)f$.
Dla $f\in\mathcal D(H)$ ilorazy różnicowe zbieżne są
w normie $L^2$ na mocy nierówności
$|(e^{-ih\lambda/\hbar}-1)/h|\le|\lambda|/\hbar$
i twierdzenia o zbieżności zdominowanej względem
miary spektralnej $\|E_H(\dd\lambda)f\|^2$.
Dostajemy pochodną $-(i/\hbar)HU(t)f$.
Analogiczny rachunek dla $S$ jest dozwolony, ponieważ
widmo $H$ jest ograniczone od dołu. Ciągłość silna
obu rodzin pochodzi z tego samego argumentu.
\end{proof}

\section{Euklidesowy propagator jest całką względem miary Wienera}

\begin{twierdzenie}[Feynman--Kac z jawnymi stałymi fizycznymi]
\label{th:physical-FK}
Niech $V\in L^\infty(\R^d;\R)$ i $H$ będzie operatorem
z \eqref{eq:QM-H}. Dla $s>0$, $f\in L^2(\R^d)$,
po wybraniu odpowiedniego reprezentanta prawej strony,
\[
 (S(s)f)(x)=\E\left[
 f\left(x+\sqrt{\frac{\hbar}{m}}B_s\right)
 \exp\left(-\frac1\hbar\int_0^s
 V\left(x+\sqrt{\frac{\hbar}{m}}B_r\right)\,\dd r\right)
 \right]
\tag{16.3}\label{eq:physical-FK}
\]
dla prawie wszystkich $x\in\R^d$, gdzie $B$ jest
$d$-wymiarowym standardowym ruchem Browna.
Dla ograniczonej funkcji borelowskiej $f$ prawa strona
jest określona w każdym $x$.
\end{twierdzenie}
\begin{proof}
Po podzieleniu równania $\hbar\partial_s v=-Hv$ przez
$\hbar$ otrzymujemy operator generatora
$\frac{\hbar}{2m}\Delta-\frac1\hbar V$.
Proces $X_s=x+\sqrt{\hbar/m}B_s$ ma generator
$\frac{\hbar}{2m}\Delta$.
Wzór Feynmana--Kaca z rozdziału~\ref{ch:path-integrals},
z potencjałem $V/\hbar$, daje więc operator $T(s)$ po
prawej stronie \eqref{eq:physical-FK}. Wykażemy, że jest
on dokładnie $S(s)$, nie zakładając istnienia klasycznego
rozwiązania PDE.

Niech $A=(\hbar/(2m))\Delta$ w $L^2$, $B=-V/\hbar$
(będzie to ograniczony operator mnożenia). Ważona
półgrupa $T(s)$ spełnia równanie całkowe Duhamela
\[
 T(s)f=e^{sA}f+\int_0^s e^{rA}B\,T(s-r)f\,\dd r.
\tag{16.4}\label{eq:FK-duhamel}
\]
Wyprowadźmy tę tożsamość bez odwoływania się do
formalnego różniczkowania półgrupy. Dla każdej trajektorii
połóżmy $g(r)=-V(X_r)/\hbar$; symbol $B=-V/\hbar$
użyty w równaniu Duhamela oznacza natomiast operator mnożenia.
Ponieważ $g$ jest ograniczona i mierzalna, z fundamentalnego
twierdzenia rachunku całkowego dla funkcji absolutnie
ciągłych otrzymujemy
\[
 e^{\int_0^s g(r)\,\dd r}
 =1+\int_0^s g(r)e^{\int_r^s g(u)\,\dd u}\,\dd r.
\]
Mnożymy przez $f(X_s)$ i bierzemy wartość oczekiwaną.
Po warunkowaniu względem $\calF_r$ czynnik pochodzący
z odcinka $[r,s]$ ma warunkową średnią
$(T(s-r)f)(X_r)$; jest to zastosowanie własności Markowa
do całej przyszłej trajektorii. Otrzymujemy
\eqref{eq:FK-duhamel} najpierw dla ograniczonych
funkcji borelowskich $f$ należących do $L^2$.
Tonelli lub Fubini wolno stosować, ponieważ
$|g|\le\|V\|_\infty/\hbar$ i waga jest ograniczona przez
$e^{s\|V\|_\infty/\hbar}$.
Wreszcie zachodzi punktowe oszacowanie
$|T(s)f|\le e^{s\|V\|_\infty/\hbar}e^{sA}|f|$.
Półgrupa cieplna jest kontrakcją w $L^2$, a więc
$\|T(s)f\|_2\le e^{s\|V\|_\infty/\hbar}\|f\|_2$.
Daje to ciągłe rozszerzenie operatorów $T(s)$ i
równania Duhamela na całe $L^2$.

Na drugiej stronie wzoru \eqref{eq:unitary-euclidean}
operator $S(s)=e^{s(A+B)}$ spełnia dokładnie to samo
równanie Duhamela dla ograniczonego zaburzenia $B$.
Jeżeli $D(s)=T(s)-S(s)$, to
$\|D(s)f\|_2\le\|B\|\int_0^s\|D(r)f\|_2\,\dd r$.
Lemat Gronwalla daje $D(s)f=0$. Wzór obowiązuje
najpierw na gęstym zbiorze ograniczonych funkcji
całkowalnych, a potem na całym $L^2$.
\end{proof}

\begin{uwaga}[Normalizacja fizyczna]
W rozdziale~\ref{ch:path-integrals} rozważaliśmy
$H=-\frac12\Delta+V$ i czas euklidesowy $t$.
Teraz pojawiają się $\hbar$ i $m$: wariancja każdej
współrzędnej przyrostu procesu $X$ wynosi
$(\hbar/m)s$, a potencjał występuje w wykładniku
z czynnikiem $1/\hbar$. Pominięcie tych stałych
prowadziłoby do błędnej skali propagatora.
\end{uwaga}

\section{Obrót Wicka w sensie twierdzenia spektralnego}

Formalna zamiana $s=it$ bywa nazywana \emph{obrotem Wicka}.
Właściwy matematyczny argument nie polega na
podstawieniu liczby zespolonej do losowej trajektorii.
Jest nim przedłużenie operatorów
$e^{-sH/\hbar}$ na prawą półpłaszczyznę
zespolonego parametru.

\begin{twierdzenie}[Przejście od półgrupy cieplnej do propagatora]
\label{th:wick-spectral}
Niech $H$ będzie samosprzężonym operatorem w zespolonej
przestrzeni Hilberta $\calH$, a jego widmo niech będzie
ograniczone od dołu. Dla $\operatorname{Re}z>0$ określmy
$S(z)=e^{-zH/\hbar}$ przez rachunek spektralny.
Wówczas $z\mapsto S(z)$ jest holomorficzne w prawej
półpłaszczyźnie w normie operatorowej lokalnie, a dla
każdego $t\in\R$ i $f\in\calH$ zachodzi
\[
 \lim_{\varepsilon\downarrow0}
  S(\varepsilon+it)f=e^{-itH/\hbar}f
  \quad\text{w normie }\calH.
\tag{16.5}\label{eq:wick-strong}
\]
Nie wynika stąd samo przez się zbieżność punktowa jąder
ani istnienie zespolonej miary na przestrzeni trajektorii.
\end{twierdzenie}
\begin{proof}
Niech $H\ge-MI$. Dla $\varepsilon\in(0,1]$ i
$\lambda\ge-M$ mamy $|e^{-\varepsilon\lambda/\hbar}|
\le e^{M/\hbar}$. Z rachunku spektralnego
\[
\begin{split}
&\|S(\varepsilon+it)f-e^{-itH/\hbar}f\|_{\calH}^2\\
&\quad=\int_{\R}\left|e^{-\varepsilon\lambda/\hbar}-1\right|^2
 \,\dd\langle E_H(\lambda)f,f\rangle.
\end{split}
\]
Całkowana funkcja dąży punktowo do zera i jest
jednostajnie ograniczona przez $(1+e^{M/\hbar})^2$.
Zbieżność zdominowana daje \eqref{eq:wick-strong}.

Holomorficzność wynika z tego, że na każdym zwartym
podzbiorze prawej półpłaszczyzny
$\lambda^k e^{-z\lambda/\hbar}$ jest jednostajnie
ograniczona dla $\lambda\ge-M$, dla każdego $k\ge0$.
Można więc różniczkować rachunek spektralny w normie
operatorowej. Tożsamość na brzegu jest jednak
\emph{zbieżnością silną operatorów}, nie całką po
trajektoriach z nową miarą probabilistyczną.
\end{proof}

\section{Cząstka swobodna: obliczenie propagatora}

\begin{twierdzenie}[Jądro cząstki swobodnej]
\label{th:free-schrodinger-kernel}
Niech $m,\hbar>0$, $H_0=-\hbar^2\Delta/(2m)$ w
$L^2(\R^d)$. Dla $t>0$ operator
$U_0(t)=e^{-itH_0/\hbar}$ ma, na funkcjach Schwartza,
oscylacyjną reprezentację
\[
 (U_0(t)f)(x)=\int_{\R^d}^{\mathrm{osc}}
 K_0(t;x,y)f(y)\,\dd\lambda_d(y),
\quad
 K_0(t;x,y)=\left(\frac{m}{2\pi i\hbar t}\right)^{d/2}
 \exp\left(\frac{im\|x-y\|^2}{2\hbar t}\right).
\tag{16.6}\label{eq:free-propagator}
\]
Potęga pierwiastka jest określona przez granicę z
$z=\varepsilon+it$, $\varepsilon>0$. Wyrażenie jest
granicą regularizowanych całek gaussowskich; nie jest
jądrem przejścia miary probabilistycznej.
\end{twierdzenie}
\begin{proof}
Stosujemy transformatę Fouriera z konwencją
$\widehat f(\xi)=\int e^{-i\langle x,\xi\rangle}
 f(x)\,\dd\lambda_d(x)$. Wtedy
\[
 \widehat{U_0(t)f}(\xi)
 =e^{-it\hbar\|\xi\|^2/(2m)}\widehat f(\xi).
\]
Dla parametru $z$ z $\operatorname{Re}z>0$
transformata odwrotna mnożnika
$e^{-z\hbar\|\xi\|^2/(2m)}$ jest klasyczną gęstością
Gaussa
\[
 q_z(x,y)=\left(\frac{m}{2\pi\hbar z}\right)^{d/2}
 \exp\left(-\frac{m\|x-y\|^2}{2\hbar z}\right).
\]
Stosujemy $q_{\varepsilon+it}$ do $f$ i dopuszczamy
$\varepsilon\downarrow0$. Twierdzenie~\ref{th:wick-spectral}
gwarantuje zbieżność w $L^2$. Dla funkcji Schwartza
można również przejść do granicy w odwróconej
transformacie Fouriera w sensie rozkładów temperowanych.
Zastępując $z$ przez $it$ w jądrach, otrzymujemy
\eqref{eq:free-propagator} ze wskazanym wyborem gałęzi.
Ponieważ $f\in\calS(\R^d)\subset L^1(\R^d)$,
sam zapis $\int K_0(t;x,y)f(y)\,\dd\lambda_d(y)$
może być odczytany jako zwykła całka absolutnie zbieżna
względem $y$. Dopisek $\mathrm{osc}$ odnosi się do
charakteru jądra i do rozszerzenia reprezentacji na
szersze klasy funkcji, dla których całka bezwzględna
nie musi istnieć.
\end{proof}

\begin{uwaga}[Dlaczego zwykła miara zespolona zawodzi już w przykładzie swobodnym?]
Dla $t>0$ moduł jądra \eqref{eq:free-propagator}
nie zależy od $y$:
$|K_0(t;x,y)|=(m/(2\pi\hbar t))^{d/2}$.
Dlatego $\int_{\R^d}|K_0(t;x,y)|\,\dd\lambda_d(y)=\infty$.
Nie można więc potraktować tego jądra jako gęstości
\emph{skończonej zespolonej miary przejścia}.
Tym bardziej nie wynika z niego istnienie przeliczalnie
addytywnej miary fazowej o skończonej wariacji na
całej przestrzeni trajektorii.
\end{uwaga}

\section{Dyskretyzacja czasu: od działania klasycznego do iloczynów operatorów}

Dla gładkiej trajektorii $x:[0,T]\to\R^d$ klasyczne
działanie cząstki w potencjale $V$ jest równe
$S[x]=\int_0^T(\frac m2\|\dot x(s)\|^2-V(x(s)))\,\dd s$.
Typowa trajektoria Brownowska nie ma jednak pochodnej:
ten wzór \emph{nie jest} funkcjonałem określonym prawie
wszędzie względem miary Wienera. Przybliżamy działanie
dla łamanych o wierzchołkach $x_0,\ldots,x_n$.

\begin{twierdzenie}[Iloczyny Trottera dla ograniczonego potencjału]
\label{th:QM-trotter}
Niech $H_0$ będzie jak w Twierdzeniu~\ref{th:free-schrodinger-kernel},
a $V\in L^\infty(\R^d;\R)$. Dla każdego $t\in\R$
i $f\in L^2(\R^d)$ zachodzi
\[
 e^{-it(H_0+V)/\hbar}f
 =\lim_{n\to\infty}
 \left(e^{-itV/(n\hbar)}e^{-itH_0/(n\hbar)}\right)^n f
 \quad\text{w }L^2(\R^d).
\tag{16.7}\label{eq:trotter-quantum}
\]
Jeżeli dodatkowo $V$ jest gładkie i wszystkie jego
pochodne są ograniczone, iloczyny można na gładkich
szybko malejących funkcjach zapisywać przez iterowane
całki oscylacyjne o dyskretnym działaniu
\[
 S_n(x_0,\ldots,x_n)
 =\sum_{k=0}^{n-1}\left(
    \frac{m\|x_{k+1}-x_k\|^2}{2\Delta}
    -\Delta V(x_k)\right),\quad\Delta=t/n,
\tag{16.8}\label{eq:discrete-action}
\]
przy czym $t>0$, $x_0=x$ i $x_n=y$ są końcami
trajektorii. Zbieżność w \eqref{eq:trotter-quantum}
nie oznacza istnienia całki Lebesgue'a z
$e^{iS[x]/\hbar}$ po nieskończonej przestrzeni trajektorii.
\end{twierdzenie}
\begin{proof}
Podajemy argument ograniczonego zaburzenia.
Niech $A=-iH_0/\hbar$ i $B=-iV/\hbar$; wtedy $A$
jest generatorem grupy unitarnej, a $B$ jest
ograniczonym operatorem antysamiosprzężonym.
Niech $U(t)=e^{t(A+B)}$ i $U_A(t)=e^{tA}$.
Wzór Duhamela dla ograniczonego zaburzenia daje
\[
 U(h)f=U_A(h)f+\int_0^h U_A(h-s)BU(s)f\,\dd s.
\]
Dla $h\downarrow0$ z ciągłości silnej $U,U_A$
wynika, że całka po prawej jest równa
$hBf+o_f(h)$. Ponadto
$e^{hB}U_A(h)f=U_A(h)f+hBf+o_f(h)$,
gdyż $B$ jest ograniczony. Mamy zatem
$h^{-1}(e^{hB}U_A(h)-U(h))f\to0$ dla każdego $f$.
Normy operatorów
$h^{-1}(e^{hB}U_A(h)-U(h))$ są jednakowo
ograniczone dla małych dodatnich $h$ na mocy
wzoru Duhamela i nierówności
$\|e^{hB}-I\|\le h\|B\|e^{h\|B\|}$.
Zbieżność silna jest wobec tego jednostajna na
zwartych zbiorach wektorów. Dla ustalonego $t>0$
zbiór $\{U(s)f:0\le s\le t\}$ jest zwarty w $L^2$.

Korzystamy z tożsamości teleskopowej, przy
$h=t/n$ i $Q_h=e^{hB}U_A(h)$:
\[
 Q_h^n f-U(t)f
 =\sum_{k=0}^{n-1} Q_h^{n-1-k}
   (Q_h-U(h))\,U(kh)f.
\]
Każdy czynnik $Q_h$ i $U(h)$ jest unitarny, dlatego
norma prawej strony jest najwyżej
$n\sup_{0\le s\le t}\|(Q_h-U(h))U(s)f\|=n\,o_f(h)=o(1)$.
Dla $t<0$ postępujemy tak samo po zamianie znaku
parametru albo stosujemy sprzężenia. Dowodzi to
\eqref{eq:trotter-quantum}.

Dla gładkiego ograniczonego $V$ każdy czynnik
$e^{-itV/(n\hbar)}$ stojący po lewej stronie czynnika
swobodnego mnoży przez fazę
$e^{-i\Delta V(x)/\hbar}$, natomiast
$e^{-itH_0/(n\hbar)}$ ma jądro z
\eqref{eq:free-propagator}. Ich iterowane składanie daje
sumę faz kinetycznych i potencjalnych
$S_n/\hbar$. Czynnik normalizacyjny
$(m/(2\pi i\hbar\Delta))^{nd/2}$ to produkt
prefaktorów jąder swobodnych. Wszystkie całki
odczytujemy jako całki oscylacyjne lub granice
regularizowane; wcześniejszy argument zapewnia
zbieżność \emph{operatorów}, a nie miar fazowych.
\end{proof}

\begin{przyklad}[Potencjał stały]
Jeżeli $V(x)\equiv v_0\in\R$, to
$[H_0,V]=0$ i $U(t)=e^{-itv_0/\hbar}U_0(t)$.
Propagator zmienia się tylko o wspólną fazę:
$K_V(t;x,y)=e^{-itv_0/\hbar}K_0(t;x,y)$.
Jest to dokładna wersja formalnego wzoru
$e^{iS[x]/\hbar}$, w którym do działania
kinetycznego dodano wyraz $-v_0T$.
\end{przyklad}

\section{Oscylator harmoniczny: przykład, który można policzyć dokładnie}

Potencjał $V(x)=\frac12 m\omega^2x^2$, $m,\omega>0$,
nie jest ograniczony. Nie podpada więc pod hipotezy
Twierdzenia~\ref{th:QM-trotter}, ale operator
$H_{\mathrm{osc}}=-\frac{\hbar^2}{2m}\frac{\dd^2}{\dd x^2}
+\frac12m\omega^2x^2$ jest samosprzężony i
ograniczony od dołu. Możemy nadal rozważyć jego
półgrupę cieplną i propagator.

\begin{twierdzenie}[Jądro Mehlera w czasie euklidesowym]
\label{th:Mehler-euclidean}
Dla $s>0$, $x,y\in\R$ jądro operatora
$e^{-sH_{\mathrm{osc}}/\hbar}$ jest równe
\[
\begin{split}
 K_s^{\mathrm E}(x,y)
  ={}&\sqrt{\frac{m\omega}{2\pi\hbar\sinh(\omega s)}}\\
 &\cdot\exp\left[-\frac{m\omega}{2\hbar\sinh(\omega s)}
 \bigl((x^2+y^2)\cosh(\omega s)-2xy\bigr)\right].
\end{split}
\tag{16.9}\label{eq:Mehler-E}
\]
Dla ograniczonej ciągłej $f$ prawa strona określa
naturalne rozszerzenie półgrupy, nawet gdy $f\notin L^2$.
W tym sensie jądro można zapisać jako
\[
\begin{split}
 (e^{-sH_{\mathrm{osc}}/\hbar}f)(x)
 =\E_x\Big[&f(X_s)\exp\Big(
  -\frac{m\omega^2}{2\hbar}\int_0^sX_r^2\,\dd r
 \Big)\Big],\\
 &X_r=x+\sqrt{\hbar/m}B_r.
\end{split}
\tag{16.10}\label{eq:Mehler-FK}
\]
\end{twierdzenie}
\begin{proof}
Najpierw sprawdzamy wzór \eqref{eq:Mehler-E}
bez odwoływania się do całek formalnych.
Szukamy dodatniego jądra w postaci
$K_s(x,y)=A_s\exp[-a_s(x^2+y^2)+b_sxy]$.
Podstawiając do równania
\[
 \partial_sK=\frac{\hbar}{2m}\partial_{xx}K
    -\frac{m\omega^2}{2\hbar}x^2K
\]
i porównując współczynniki wyrazów $x^2$, $xy$
oraz stałego, dostajemy układ
\[
 a_s'= -\frac{2\hbar}{m}a_s^2+
       \frac{m\omega^2}{2\hbar},\qquad
 b_s'=-\frac{2\hbar}{m}a_sb_s,\qquad
 \frac{A_s'}{A_s}=-\frac{\hbar}{m}a_s,
\tag{16.11}\label{eq:Mehler-odes}
\]
oraz warunek zgodności
$b_s^2=4a_s^2-(m\omega/\hbar)^2$.
Rozwiązanie spełniające asymptotykę jądra Gaussa
przy $s\downarrow0$ jest równe
\[
 a_s=\frac{m\omega}{2\hbar}\coth(\omega s),\quad
 b_s=\frac{m\omega}{\hbar\sinh(\omega s)},\quad
 A_s=\sqrt{\frac{m\omega}{2\pi\hbar\sinh(\omega s)}}.
\]
Rzeczywiście $\coth' r=-1/\sinh^2r$ oraz
$\coth^2r-1=1/\sinh^2r$ weryfikują
\eqref{eq:Mehler-odes}. Gdy $s\downarrow0$,
$\sinh(\omega s)\sim\omega s$ i
$\cosh(\omega s)\sim1$, więc $K_s(x,y)$ jest
asymptotycznie jądrem
$\sqrt{m/(2\pi\hbar s)}
\exp[-m(x-y)^2/(2\hbar s)]$.
Z własności przybliżonej jedności wynika
$\int K_s(x,y)f(y)\,\dd y\to f(x)$
dla ciągłych funkcji o zwartym nośniku.

Sprawdźmy własność półgrupową, której nie wolno
wnioskować tylko z podstawienia do PDE. Przy oznaczeniach
$a_s,b_s,A_s$ wprowadzonych powyżej z całki Gaussa wynika
\[
\begin{split}
&\int_{\R}K_s(x,z)K_t(z,y)\,\dd z\\
&=A_sA_t\sqrt{\frac{\pi}{a_s+a_t}}
\exp\!\left[-\left(a_s-\frac{b_s^2}{4(a_s+a_t)}\right)x^2
-\left(a_t-\frac{b_t^2}{4(a_s+a_t)}\right)y^2
+\frac{b_sb_t}{2(a_s+a_t)}xy\right].
\end{split}
\]
Wstawiając wzory na $a_s,b_s,A_s$ i korzystając z
$\sinh(\omega(s+t))=\sinh(\omega s)\cosh(\omega t)
+\cosh(\omega s)\sinh(\omega t)$, sprawdzamy, że
prefaktor wynosi $A_{s+t}$, oba współczynniki przy
$x^2,y^2$ są równe $-a_{s+t}$, a współczynnik przy
$xy$ jest równy $b_{s+t}$. Otrzymujemy więc dokładnie
$K_{s+t}(x,y)$.
Dla $s>0$ całkowanie Gaussa daje też
$\int K_s(x,y)\,\dd y\le1$; symetria jądra i
oszacowanie Schura zapewniają
$\|K_sf\|_2\le\|f\|_2$.
Własność przybliżonej jedności daje silną ciągłość
w zerze na gęstym $C_c^\infty$, a oszacowanie normy
rozszerza ją na całe $L^2$. Samą identyfikację
generatora z domknięciem operatora oscylatora
uzasadnimy poniżej przez jednoznaczność równania
w $L^2$; nie zakładamy jej na podstawie samej
zgodności formalnych pochodnych.

Do reprezentacji probabilistycznej nie potrzebujemy
ogólnego twierdzenia o monotonicznej zbieżności form.
Dla $f\in C_c^\infty(\R)$, $f\ge0$, połóżmy
$u(s,x)=\int_{\R}K_s^{\mathrm E}(x,y)f(y)\,\dd y$.
Obliczając całkę Gaussa, dostajemy
\[
 \int_{\R} K_s^{\mathrm E}(x,y)\,\dd y
 =\frac{\exp[-(m\omega/(2\hbar))x^2\tanh(\omega s)]}
        {\sqrt{\cosh(\omega s)}}\le1.
\]
Zatem $0\le u(s,x)\le\|f\|_\infty$.
Funkcja $u$ spełnia klasyczne równanie
$\partial_su=(\hbar/(2m))u_{xx}-(m\omega^2/(2\hbar))x^2u$
przy $s>0$ oraz $u(s,x)\to f(x)$ lokalnie jednostajnie,
gdy $s\downarrow0$.

Niech $X_r=x+\sqrt{\hbar/m}B_r$. Dla
$0<\varepsilon<s$ i $R>|x|$ stosujemy wzór Itô do
\[
 Z_r=\exp\left(-\frac{m\omega^2}{2\hbar}
 \int_0^rX_v^2\,\dd v\right)u(s-r,X_r),
 \qquad 0\le r\le(s-\varepsilon)\wedge\tau_R,
\]
gdzie $\tau_R=\inf\{r\ge0:|X_r|\ge R\}$.
Dryf znika dzięki równaniu PDE, a po zatrzymaniu
współczynnik stochastyczny jest ograniczony na
odpowiednim zwartym zbiorze czasu i przestrzeni,
więc $Z$ jest martyngałem. Mamy $0\le Z\le\|f\|_\infty$.
Zatem $\E Z_{(s-\varepsilon)\wedge\tau_R}=u(s,x)$.
Dopuszczamy najpierw $R\to\infty$: ciągłość Brownowskich
trajektorii daje $\tau_R\to\infty$ prawie na pewno,
a zbieżność zdominowana uzasadnia przejście do granicy.
Następnie $\varepsilon\downarrow0$; ciągłość $X$ i
lokalnie jednostajna zbieżność $u(\varepsilon,\cdot)\to f$
dają właśnie \eqref{eq:Mehler-FK}.
Dla $f$ dowolnej ograniczonej ciągłej stosujemy obcięcia
i zbieżność zdominowaną, a dla funkcji zmieniających znak
rozkład na części dodatnią i ujemną.

Aby zidentyfikować otrzymaną półgrupę z
$e^{-sH_{\mathrm{osc}}/\hbar}$ w $L^2$, wykorzystujemy
jednoznaczność rozwiązań równania cieplnego z
potencjałem nieujemnym w klasie $L^2$: dla różnicy
dwóch rozwiązań z zerową daną początkową iloczyn skalarny
z rozwiązaniem i całkowanie przez części dają
\[
 \frac{\dd}{\dd s}\|v(s)\|_2^2
 =-\frac{\hbar}{m}\|\partial_xv(s)\|_2^2
  -\frac{m\omega^2}{\hbar}\|xv(s)\|_2^2\le0.
\]
Wyjaśnijmy dokładniej przejście poza funkcje o zwartym
nośniku. Wybierzmy $\chi\in C_c^\infty(\R)$, $0\le\chi\le1$,
równą $1$ na $[-1,1]$, i połóżmy $\chi_R(x)=\chi(x/R)$.
Mnożąc równanie dla różnicy $v$ przez $\chi_R^2\bar v$,
całkując i biorąc części rzeczywiste, otrzymujemy po
nierówności Cauchy'ego--Schwarza
\[
\frac{\dd}{\dd s}\int\chi_R^2|v|^2\,\dd x
\le \frac{C}{R^2}\int_{\R}|v|^2\,\dd x.
\]
Istotnie, wyraz mieszany zawiera
$2\chi_R\chi_R'\bar v\,\partial_xv$ i można go
wchłonąć w ujemny wyraz
$-\frac{\hbar}{m}\int\chi_R^2|\partial_xv|^2$,
pozostawiając $C\int|\chi_R'|^2|v|^2$;
składnik potencjału jest niedodatni. Na przedziale
$[\delta,s]$ rachunek jest klasyczny dla rozwiązań
gładkich; w przestrzeni formowej otrzymuje się go przez
aproksymację w normie formy. Dopuszczamy następnie
$R\to\infty$ i $\delta\downarrow0$, korzystając z
ciągłości w $L^2$ i ograniczenia normy na zwartych
przedziałach czasu. Dostajemy
$\|v(s)\|_2^2\le\|v(0)\|_2^2=0$.
Półgrupa spektralna i operator z jądrem $K_s^{\mathrm E}$
dają rozwiązania z tymi samymi danymi, więc są identyczne.
Nie wykorzystaliśmy tutaj nieudowodnionej jednoznaczności
klasycznej dla dowolnych nieograniczonych potencjałów.
\end{proof}

\begin{wniosek}[Propagator oscylatora w czasie rzeczywistym]
\label{cor:Mehler-real}
Niech $0<t<\pi/\omega$. Wtedy jądro propagatora
$e^{-itH_{\mathrm{osc}}/\hbar}$ ma oscylacyjną postać
\[
\begin{split}
 K_t^{\mathrm R}(x,y)
  ={}&\sqrt{\frac{m\omega}{2\pi i\hbar\sin(\omega t)}}\\
 &\cdot\exp\left[\frac{im\omega}{2\hbar\sin(\omega t)}
 \bigl((x^2+y^2)\cos(\omega t)-2xy\bigr)\right].
\end{split}
\tag{16.12}\label{eq:Mehler-R}
\]
Dla dłuższych czasów konieczny jest wybór fazy
zgodnej z ciągłością grupy unitarnej; w chwilach
$\sin(\omega t)=0$ propagator nie ma zwykłego jądra
funkcyjnego tej postaci.
\end{wniosek}
\begin{proof}
Stosujemy Twierdzenie~\ref{th:wick-spectral} do
$H_{\mathrm{osc}}$ i dla $z=\varepsilon+it$ korzystamy
z analitycznej kontynuacji wzoru \eqref{eq:Mehler-E}.
Ponieważ $\sinh(i\omega t)=i\sin(\omega t)$ i
$\cosh(i\omega t)=\cos(\omega t)$, otrzymujemy
\eqref{eq:Mehler-R} jako granicę oscylacyjnych jąder
przy $\varepsilon\downarrow0$. W przedziale
$0<t<\pi/\omega$ gałąź pierwiastka jest wybrana
jednoznacznie przez zbieżność od $t=0+$.
Granica operatorów jest silna w $L^2$, co uzasadnia
propagator, bez wymagania istnienia miary zespolonej
na przestrzeni trajektorii.
\end{proof}

\begin{zadanie}[Przypadek zerowej częstości]
Sprawdzić, że gdy $\omega\downarrow0$, jądro
\eqref{eq:Mehler-E} zbiega do jądra cząstki swobodnej
w czasie euklidesowym.

\emph{Rozwiązanie.} Korzystamy z
$\sinh(\omega s)\sim\omega s$,
$\cosh(\omega s)=1+O(\omega^2s^2)$.
Prefaktor dąży do $\sqrt{m/(2\pi\hbar s)}$,
a wykładnik do $-m(x-y)^2/(2\hbar s)$.
Zatem $K_s^{\mathrm E}(x,y)\to q_s(x,y)$.
W tym rachunku zanik potencjału prowadzi dokładnie do
półgrupy cieplnej cząstki swobodnej.
\end{zadanie}


\section{Ślad półgrupy oscylatora i funkcja podziału}

\begin{zadanie}[Gaussowski rachunek śladu]
Dla $H_{\mathrm{osc}}$ z
Twierdzenia~\ref{th:Mehler-euclidean} obliczyć
$\int_{\R}K_s^{\mathrm E}(x,x)\,\dd x$ i porównać
z sumą po poziomach energetycznych.

\emph{Rozwiązanie.} Podstawiając $y=x$ do
\eqref{eq:Mehler-E} i korzystając z
$(\cosh r-1)/\sinh r=\tanh(r/2)$, mamy
\[
 K_s^{\mathrm E}(x,x)=
 \sqrt{\frac{m\omega}{2\pi\hbar\sinh(\omega s)}}
 \exp\left(-\frac{m\omega}{\hbar}
     \tanh(\omega s/2)x^2\right).
\]
Całka Gaussa daje
\[
 \int_{\R}K_s^{\mathrm E}(x,x)\,\dd x
 =\frac1{2\sinh(\omega s/2)}.
\]
Aby nie zakładać śladowości bez uzasadnienia,
zauważmy, że dodatnie jądro jest symetryczne i spełnia
własność półgrupową. Wobec tego
\[
\int_{\R^2}|K_{s/2}^{\mathrm E}(x,y)|^2\,
\dd\lambda_2(x,y)=\int_{\R}K_s^{\mathrm E}(x,x)\,
\dd\lambda_1(x)<\infty.
\]
Operator o jądrze $K_{s/2}^{\mathrm E}$ jest zatem
Hilberta--Schmidta, a jego kwadrat, równy
$e^{-sH_{\mathrm{osc}}/\hbar}$, jest dodatnim
operatorem śladowym. To uzasadnia utożsamienie
całki przekątnej jądra z jego śladem.
Znane widmo oscylatora,
$E_n=\hbar\omega(n+\frac12)$, $n\ge0$,
daje niezależne sprawdzenie:
\[
 \operatorname{Tr}(e^{-sH_{\mathrm{osc}}/\hbar})
 =\sum_{n=0}^\infty e^{-s\omega(n+1/2)}
 =\frac1{2\sinh(\omega s/2)}.
\]
W ten sposób jądro otrzymane z całki po trajektoriach
koduje wszystkie poziomy energetyczne. Śladowość
uzasadniliśmy powyżej; wzór na widmo oscylatora
pozostaje przywołanym wynikiem analizy operatorowej,
nie jest wyprowadzany z samego rachunku całki Gaussa.
W języku mechaniki statystycznej parametr $s$ odpowiada
odwrotnej temperaturze przez $s=\beta\hbar$;
funkcja podziału wynosi wtedy
$Z(\beta)=\operatorname{Tr}(e^{-\beta H_{\mathrm{osc}}})$.
\end{zadanie}

\section{Co jest dowiedzione, a co pozostaje formalną analogią?}

Nasze rozróżnienie nie jest jedynie słowne. W czasie
euklidesowym wzór \eqref{eq:physical-FK} jest zwykłą
wartością oczekiwaną po przestrzeni trajektorii, o ile
potencjał spełnia przyjęte założenia. W czasie rzeczywistym
propagator istnieje jako grupa unitarna z twierdzenia
spektralnego. Dla cząstki swobodnej i oscylatora
harmonicznego znamy nawet jawne jądra oscylacyjne.
Natomiast formalny zapis
\[
 \int_{x(0)=x}^{x(T)=y}
  \exp\left(\frac{i}{\hbar}S[x]\right)\,\mathcal D x
\]
\emph{nie został} zinterpretowany jako całka Lebesgue'a
względem przeliczalnie addytywnej miary zespolonej na
przestrzeni wszystkich trajektorii.
Wynik \eqref{eq:trotter-quantum} daje ścisły sens
aproksymacjom \emph{operatorowym} przez skończone iloczyny,
a Twierdzenie~\ref{th:wick-spectral} daje ścisły sens
przejściu między dodatnią półgrupą a unitarną ewolucją.
Nie należy utożsamiać tych faktów z istnieniem
ogólnej miary Feynmana.

\begin{center}
\begin{tikzpicture}[>=Latex, node distance=8mm and 12mm]
\node[draw,rounded corners,align=center,text width=5.5cm,
 minimum height=1.05cm] (a)
 {\textbf{Czas euklidesowy $s>0$}\\
  $e^{-sH/\hbar}$, miara Wienera, dodatnie wagi};
\node[draw,rounded corners,align=center,text width=5.5cm,
 minimum height=1.05cm,right=of a] (b)
 {\textbf{Czas rzeczywisty $t\in\R$}\\
  $e^{-itH/\hbar}$, ewolucja unitarna};
\draw[->,thick,blue!70!black] (a)--(b);
\node[draw,rounded corners,align=center,text width=5.5cm,
 minimum height=0.9cm,below=14mm of a] (c)
 {Całka po trajektoriach\\
  jako wartość oczekiwana};
\node[draw,rounded corners,align=center,text width=5.5cm,
 minimum height=0.9cm,below=14mm of b] (d)
 {Oscylacyjne całki skończeniewymiarowe\\
  i granice iloczynów operatorów};
\draw[->,thick] (a)--(c);
\draw[->,thick] (b)--(d);
\draw[->,densely dashed] (c)--(d);
\end{tikzpicture}
\par\smallskip\small
Strzałka górna oznacza silną granicę spektralną; strzałka
przerywana przypomina, że przejście do całek oscylacyjnych
nie jest zwykłą zmianą miary na przestrzeni trajektorii.
\end{center}

\paragraph{Zakończenie tomu III.}
Droga od izometrii Wienera, przez całkę Itô, martyngały,
SDE, procesy Lévy'ego i półmartyngały, prowadzi do
probabilistycznych reprezentacji półgrup. Całki po
trajektoriach są ścisłym narzędziem analizy; ich związek
z całkami Feynmana i mechaniką kwantową jest potężny,
ale wymaga kontrolowania samosprzężoności operatorów,
regularności potencjałów i rodzaju używanej zbieżności.

